\documentclass[10pt,a4paper]{amsart}
\usepackage[margin=19mm]{geometry}
\usepackage{amsmath,amssymb,mathtools}
\usepackage[scr=rsfs,cal=euler]{mathalfa}
\usepackage{xfrac}
\usepackage[unicode,hidelinks,bookmarks=false]{hyperref}
\newcommand{\ie}{i.e.\ }
\newcommand{\cf}{cf.\ }
\newcommand{\resp}{respectively\ }
\newcommand{\Subsection}{\subsection}
\providecommand{\nobreakdash}{\nobreak\hbox{-}}
\setlength{\emergencystretch}{2em}
\multlinegap0pt
\usepackage{mathtools}
\usepackage{cleveref}
\newtheorem{theorem}{Theorem}[section]
\newtheorem{lemma}[theorem]{Lemma}
\newtheorem{proposition}[theorem]{Proposition}
\newtheorem{conjecture}[theorem]{Conjecture}
\newtheorem*{remark}{Remark}
\newcommand{\E}{\mathbb E}
\newcommand{\Prob}{\mathbb P}
\newcommand{\cP}{\mathcal P}
\newcommand{\one}{\mathbf 1}
\DeclareMathOperator{\avdeg}{\overline d}
\setlength{\emergencystretch}{2em}
\begin{document}
\title[Nearly Spanning Regular Subgraphs]{Nearly Spanning Regular Subgraphs}
\author{Varun Sivashankar}
\date{}
\maketitle
\noindent\emph{Source and licence.} Nearly Spanning Regular Subgraphs. Version arXiv:2609.19777v1 (CC BY 4.0). This material is distributed under the Creative Commons Attribution 4.0 International licence, http://creativecommons.org/licenses/by/4.0/, as recorded in the arXiv OAI licence metadata for this exact version and shown on the abstract page https://arxiv.org/abs/2609.19777v1. Full rights notice: licenses/arxiv-2609.19777-CC-BY-4.0.txt.\par\smallskip
\noindent\textit{Editorial scope.} Counted unit: the original \texttt{theorem} environment \texttt{thm:k2} at source lines 389--394 together with its complete proof at lines 503--516; the delivered document keeps source lines 387--517 so that every referenced label and every citation resolves inside it. Native editable source is the paper's own concatenated TeX; the AMS wrapper is editorial formatting only. No publisher class, upstream script or unrelated artwork is redistributed.
\section{The sharp degree-two case}\label{sec:k2}


\begin{theorem}\label{thm:k2}
For every odd integer $r\ge3$,
\[
\delta_2(r)=\frac1{r^2-3}.
\]
\end{theorem}


For even $r$, Petersen's $2$-factorization theorem~\cite{Petersen1891} gives $\delta_2(r)=0$. The case $r=3$ of \cref{thm:k2} is due to Choi et al.~\cite{CKKPW2019}.

We follow the approach of Choi et al.: delete the bridges and bound the number of vertices omitted in each remaining component. Their argument does not generalize to odd $r > 3$. We extend their approach by adding loops to represent omitted vertices and applying the degree-two case of Edmonds' theorem. The resulting estimate also proves the bound proposed by Choi et al. in terms of bridges and degree deficiency; see \cref{prop:k2-bridges}.

\subsection{The local estimate}

Throughout this section, let $r\ge3$ be odd. Multigraphs may have loops, each contributing two to the degree. For a multigraph $B$ with $\Delta(B)\le r$, define its degree deficiency by
\[
D_r(B)=\sum_{v\in V(B)}(r-d_B(v)).
\]
For $r=3$, the following estimate appears in the proof of Choi et al.'s Theorem~2.3~\cite{CKKPW2019}.

\begin{lemma}\label[lemma]{lem:k2-local}
Let $r\ge3$ be odd, and let $B$ be a connected bridgeless multigraph with $\Delta(B)\le r$. Then $B$ has a $2$-regular subgraph omitting at most
\[
\left\lfloor\frac{D_r(B)}r\right\rfloor
\]
vertices.
\end{lemma}

To prove the lemma, we add one loop at each vertex of $B$. A $2$-factor of the enlarged graph consists of a $2$-regular subgraph of $B$ together with an added loop at each omitted vertex. Thus we seek a $2$-factor using few added loops.

We use the following consequence of Edmonds' theorem~\cite{Edmonds1965}; see Kobayashi~\cite[Section~2]{Kobayashi2022} for a formulation allowing loops. For a multigraph $J$ and $S\subseteq V(J)$, let $\partial_J S$ denote the set of edges with exactly one endpoint in $S$.

\begin{lemma}\label[lemma]{lem:edmonds-two-factor}
Let $J$ be a multigraph and assign a number $x_e\in[0,1]$ to each edge $e\in E(J)$. Suppose that the weights incident with each vertex sum to two, counting loops twice, and that
\begin{equation}\label{eq:two-factor-cut}
\sum_{e\in\partial_J S\setminus A}x_e
+\sum_{e\in A}(1-x_e)\ge1
\end{equation}
for every $S\subseteq V(J)$ and every subset $A\subseteq\partial_J S$ of odd size. Then there is a probability distribution on the $2$-factors $F$ of $J$ such that $\Prob(e\in E(F))=x_e$ for every edge $e$.
\end{lemma}

The cut condition has a simple interpretation. A $2$-factor uses an even number of edges across every cut, so its selected cut edges cannot equal an odd set $A$. It must therefore use an edge of $\partial_J S\setminus A$ or omit an edge of $A$. Taking expectations gives \eqref{eq:two-factor-cut}. Edmonds' theorem states that these conditions, together with the degree and weight conditions, suffice.

\begin{proof}[Proof of \cref{lem:k2-local}]
Add a loop $\ell_v$ at every vertex $v$ of $B$, obtaining a multigraph $B^+$. Set
\[
x_e=\frac2r\quad(e\in E(B)),
\qquad
x_{\ell_v}=\frac{r-d_B(v)}r\quad(v\in V(B)).
\]
These weights lie in $[0,1]$. At each vertex $v$, their sum, counting loops twice, is
\[
\frac{2d_B(v)}r+2\frac{r-d_B(v)}r=2.
\]

We check \eqref{eq:two-factor-cut}. Fix $S\subseteq V(B^+)$ and an odd set $A\subseteq\partial_{B^+}S$, and write $c=|\partial_{B^+}S|$ and $a=|A|$. Loops do not cross cuts, so every edge in this cut belongs to $B$ and has weight $2/r$. Since $a\ge1$ and $B$ has no bridge, $c\ge2$. If $a=1$, the left side of \eqref{eq:two-factor-cut} is
\[
\frac2r(c-1)+1-\frac2r
=1+\frac{2(c-2)}r\ge1.
\]
If $a\ge3$, then $c\ge a$ and $r\ge3$ give
\[
\frac2r(c-a)+\left(1-\frac2r\right)a
\ge\frac a3\ge1.
\]

Lemma~\ref{lem:edmonds-two-factor} therefore gives a random $2$-factor $F$ of $B^+$ for which the expected number of added loops is
\[
\E\bigl|\{v:\ell_v\in E(F)\}\bigr|
=\sum_{v\in V(B)}x_{\ell_v}
=\frac{D_r(B)}r.
\]
Some $2$-factor consequently uses at most $\lfloor D_r(B)/r\rfloor$ added loops. Delete the vertices carrying those loops. No other selected edge is incident with such a vertex, so every remaining vertex keeps degree two. The remaining subgraph lies in $B$ and has the required order.
\end{proof}

\subsection{Bridges and the upper bound}

Recall that $f_2(G)$ is the largest number of vertices in a $2$-regular subgraph of $G$. The local estimate yields the following bound, proposed with a ceiling in place of the floor by Choi et al.~\cite[Section~1]{CKKPW2019}.

\begin{proposition}\label[proposition]{prop:k2-bridges}
Let $r\ge3$ be odd, and let $G$ be a multigraph with $\Delta(G)\le r$, $c$ bridges and degree deficiency $d=D_r(G)=\sum_{v\in V(G)}(r-d_G(v))$. Then
\[
|V(G)|-f_2(G)\le
\max\left\{0,\left\lfloor\frac{d+c-1}{r-1}\right\rfloor\right\}.
\]
\end{proposition}

\begin{proof}
Delete all bridges of $G$, obtaining components $B_1,\ldots,B_m$. Contracting each $B_i$ to a vertex $i$ gives a forest $T$ with $c$ edges. The degree deficiency of $B_i$ is
\[
D_r(B_i)=\sum_{v\in V(B_i)}(r-d_G(v))+d_T(i),
\]
since each bridge incident with $B_i$ contributes one additional unit of deficiency after deletion.

Put $t_i=\lfloor D_r(B_i)/r\rfloor$ and $t=\sum_i t_i$. By \cref{lem:k2-local}, the union of suitable $2$-regular subgraphs of the $B_i$ omits at most $t$ vertices. If $t=0$, the assertion follows. Otherwise, let $I=\{i:t_i>0\}$, so $1\le|I|\le t$.

In the sum $\sum_{i\in I}d_T(i)$, each edge of $T$ is counted at most once except for edges of the induced forest $T[I]$, which are counted twice. Hence
\[
\sum_{i\in I}d_T(i)
\le c+e(T[I])\le c+|I|-1 \le c+t-1.
\]
Since $rt_i\le D_r(B_i)$ and each term in the sum defining $d$ is nonnegative, we obtain
\[
\begin{aligned}
rt
&\le\sum_{i\in I}D_r(B_i)
\le d+\sum_{i\in I}d_T(i)
\le d+c+t-1.
\end{aligned}
\]
Thus $t\le(d+c-1)/(r-1)$, and taking the floor proves the assertion.
\end{proof}

In particular, an $r$-regular graph with at most $r-1$ bridges has a $2$-factor, recovering the theorem of Hanson, Loten and Toft~\cite{HLT1998}. For connected graphs, van den Heuvel and Toft~\cite[Theorem~1.9]{HT2026} replace the bridge bound by the weaker requirement that at most \(r-1\) components obtained by deleting the bridges are incident with exactly one bridge.


\begin{proof}[Proof of the upper bound in \cref{thm:k2}]
Let $G$ be a connected simple $r$-regular graph on $n$ vertices, where $r\ge3$ is odd, and let $c$ be its number of bridges. If $c=0$, then \cref{lem:k2-local} gives a $2$-factor. If $c\ge1$, the bridge-count bound of O and West~\cite[Theorem~3.4]{OW2010} gives
\[
c-1\le\frac{(r-1)(n-2(r+2))}{r^2-3}.
\]
Since $D_r(G)=0$, \cref{prop:k2-bridges} now yields
\[
|V(G)|-f_2(G)
\le\left\lfloor\frac{c-1}{r-1}\right\rfloor
\le\left\lfloor\frac{n-2(r+2)}{r^2-3}\right\rfloor
<\frac{n}{r^2-3}.
\]
Applying this bound separately to the connected components proves $\delta_2(r)\le1/(r^2-3)$.
\end{proof}

\begin{thebibliography}{99}
\bibitem[CKKPW2019]{CKKPW2019} I.~Choi, R.~Kim, A.~V.~Kostochka, B.~Park and D.~B.~West, Largest $2$-regular subgraphs in $3$-regular graphs, \emph{Graphs and Combinatorics} \textbf{35} (2019), 805--813. \href{https://doi.org/10.1007/s00373-019-02021-6}{doi:10.1007/s00373-019-02021-6}.
\bibitem[Edmonds1965]{Edmonds1965} J.~Edmonds, Maximum matching and a polyhedron with $0,1$-vertices, \emph{Journal of Research of the National Bureau of Standards, Section B} \textbf{69B} (1965), 125--130. \href{https://doi.org/10.6028/jres.069B.013}{doi:10.6028/jres.069B.013}.
\bibitem[HLT1998]{HLT1998} D.~Hanson, C.~O.~M.~Loten and B.~Toft, On interval colourings of bi-regular bipartite graphs, \emph{Ars Combinatoria} \textbf{50} (1998), 23--32.
\bibitem[HT2026]{HT2026} J.~van den Heuvel and B.~Toft, $2$-Factors in Graphs, \emph{Electronic Journal of Combinatorics} \textbf{33} (2026), no.~2, P2.40. \href{https://doi.org/10.37236/14739}{doi:10.37236/14739}.
\bibitem[Kobayashi2022]{Kobayashi2022} Y.~Kobayashi, Weighted triangle-free $2$-matching problem with edge-disjoint forbidden triangles, \emph{Mathematical Programming} \textbf{192} (2022), 675--702. \href{https://doi.org/10.1007/s10107-021-01661-y}{doi:10.1007/s10107-021-01661-y}.
\bibitem[OW2010]{OW2010} S.~O and D.~B.~West, Balloons, cut-edges, matchings, and total domination in regular graphs of odd degree, \emph{Journal of Graph Theory} \textbf{64} (2010), 116--131. \href{https://doi.org/10.1002/jgt.20443}{doi:10.1002/jgt.20443}.
\bibitem[Petersen1891]{Petersen1891} J.~Petersen, Die Theorie der regul\"aren graphs, \emph{Acta Mathematica} \textbf{15} (1891), 193--220. \href{https://doi.org/10.1007/BF02392606}{doi:10.1007/BF02392606}.
\end{thebibliography}
\end{document}
